\chapter{Introdução}

O dimensionamento de equipamentos em processos químicos é uma etapa crítica no projeto de plantas industriais, envolvendo princípios fundamentais da engenharia química, como balanços de massa e energia, termodinâmica, transferência de calor e de massa e cinética química. Esses conceitos são aplicados para definir características operacionais e tamanhos de equipamentos, como trocadores de calor, colunas de destilação, reatores, bombas e compressores. O domínio das operações unitárias é essencial para garantir desempenho seguro, eficiente e otimizado em termos de custos nos processos industriais \cite{felder2005}.

Apesar de \textit{softwares} consagrados como \textit{Aspen Plus}\textregistered, \textit{HYSYS}\textregistered\ e \textit{ChemCAD}\textregistered\ serem amplamente utilizados na indústria, seu uso é limitado por custos elevados de licenciamento e pela falta de acesso em ambientes acadêmicos \cite{aspen2025, chemcad2025, alves2022}.

Essas ferramentas, embora poderosas, nem sempre priorizam o aprendizado conceitual, o que pode dificultar a compreensão dos cálculos pelos estudantes e engenheiros em formação. Há, portanto, espaço para soluções acessíveis, flexíveis e didáticas. Este trabalho propõe o desenvolvimento de um \textit{software} \textit{web} próprio e gratuito para simulação e dimensionamento de operações unitárias, integrando conhecimento teórico com programação moderna.